\documentclass{situationist-preprint}

\title{Against Higher-Order Theories}
\msid{MND-14-a}
\msfamily{MND-14}
\msdiscipline{Philosophy of Mind}
\mstype{Principal manuscript}
\mskeywords{higher-order thought; higher-order perception; misrepresentation; targetless representation; transitivity principle; phenomenal consciousness}

\begin{document}
\maketitle

\begin{abstract}
Higher-order theories hold that a mental state is conscious when, and because, the subject represents herself as being in it by a suitable higher-order state. Such theories must say what happens when the higher-order state misrepresents: when it represents the subject as being in a state of one kind while she is in a state of another, or in no state at all. We argue that consciousness is not a higher-order representation, because no higher-order theory can handle misrepresentation without positing consciousness of non-existent states. In the targetless case there are three options: nothing is conscious, the non-existent first-order state is conscious, or the higher-order state is itself the conscious state. The first contradicts the theory's own account of what it is like for the subject; the second posits a conscious state that does not exist; the third abandons the claim that a state is conscious because it is represented, and leaves the theory with an unrepresented conscious state. We show that the mismatch case reduces to the targetless case, and reply to the responses of Rosenthal, Weisberg and Brown.
\end{abstract}

\section{The problem}

On higher-order theories, what makes a mental state conscious is that its subject is aware of being in it, and this awareness is a further mental state that represents the first. The theories differ in the nature of the higher-order state: a thought, on the higher-order thought theory of \citet{Rosenthal2005}; a quasi-perceptual state of inner sense, on the theories of \citet{Armstrong1968} and \citet{Lycan1996}; a disposition to produce such a thought, on the dispositional theory of \citet{Carruthers2000}. All accept what Rosenthal calls the transitivity principle: a conscious state is one the subject is in some way aware of being in. They have received support from evidence that higher-order cortical areas contribute to awareness \citep{LauRosenthal2011}.

Since higher-order states are representations, they can misrepresent. \citet{Byrne1997} and \citet{Neander1998} pressed the question of what such theories should say about the case in which a higher-order state represents a first-order state that does not exist, and \citet{Block2011} argued that every answer is fatal. Defenders have replied that the theory was always committed to the view that what it is like for the subject is fixed by the higher-order state \citep{Rosenthal2011,Weisberg2011}, or that the theory should be restated so that the higher-order state is the phenomenally conscious state \citep{Brown2015}. We argue that the dilemma is sound and that none of these replies escapes it. MND-23 relies on this result in setting higher-order theories aside when it tests its verdicts on animal consciousness against the main theories.

\section{Framework}

\begin{definition}[Higher-order theory]\label{d:ho}
A \emph{higher-order theory} of consciousness holds that a mental state $M$ of a subject $S$ is conscious if and only if, and because, $S$ is in a suitable higher-order state $H$ that represents $S$ as being in $M$.
\end{definition}

\begin{definition}[Misrepresentation]\label{d:mis}
A higher-order state $H$ that represents $S$ as being in a state of kind $K$ is \emph{targetless} if $S$ is in no state of kind $K$, and \emph{mismatched} if $S$ is in a state of another kind $K'$ that $H$ is caused by or directed at.
\end{definition}

\begin{assumption}[Possibility of misrepresentation]\label{a:poss}
For any higher-order state of the kind a higher-order theory invokes, targetless and mismatched states of that kind are possible.
\end{assumption}

Assumption~\ref{a:poss} is accepted by defenders of the theories; Rosenthal holds that higher-order thoughts are distinct existences from their targets and can occur without them, and the empirical evidence the theories cite concerns dissociations between first-order processing and awareness.

\begin{principle}[Existence]\label{p:exist}
Only existing states are conscious. If there is something it is like for a subject at a time, then there is a state of that subject at that time in virtue of which it is so.
\end{principle}

\begin{principle}[Subjective sameness]\label{p:same}
On any higher-order theory, a subject in a targetless state $H$ is, from her own point of view, as she would be if $H$ had its target: she is aware of herself as being in $M$, and what it is like for her is the same.
\end{principle}

Principle~\ref{p:same} is not our premise but the theories' own: they explain the subjective character of a conscious state by the way the subject represents herself, and a targetless $H$ represents her in the same way \citep{Rosenthal2011}.

\section{The argument}

\begin{lemma}\label{l:three}
In a targetless case, a higher-order theory must say one of three things: (a) nothing is conscious; (b) the non-existent state $M$ is conscious; (c) $H$ is itself a conscious state, without being the target of a further higher-order state.
\end{lemma}

\begin{proof}
By Principle~\ref{p:same}, there is something it is like for the subject, so the theory must locate a conscious state or deny that there is one. The only candidates are $M$, which does not exist, and $H$, which exists. If $H$ were conscious because it is the target of a further higher-order state, the question would arise again for that state; we consider the simple case in which $H$ is not so targeted, which is the normal case on all the theories, since most conscious states are not introspected.
\end{proof}

\begin{theorem}\label{t:main}
No higher-order theory handles targetless misrepresentation without positing consciousness of non-existent states. Therefore consciousness is not a higher-order representation.
\end{theorem}

\begin{proof}
By Assumption~\ref{a:poss}, targetless cases are possible. Take the options of Lemma~\ref{l:three}. (a) contradicts Principle~\ref{p:same}: the theory says that what it is like for the subject is fixed by her higher-order representation of herself, and that representation is present. (c) contradicts Definition~\ref{d:ho}: $H$ is conscious but is not represented by any higher-order state, so a state is conscious without being the target of a higher-order representation, which is what the theory denies; the theory has become a theory on which the conscious state is the higher-order state itself, and higher-order representation is no longer what makes a state conscious. This leaves (b), on which the conscious state is a state that does not exist. That is to posit consciousness of a non-existent state, and it violates Principle~\ref{p:exist}. Since Definition~\ref{d:ho} requires one of the three, and only (b) is consistent with Definition~\ref{d:ho} and Principle~\ref{p:same}, a higher-order theory either posits consciousness of non-existent states or fails as a theory of consciousness. Positing such states is unacceptable by Principle~\ref{p:exist}. So consciousness is not a higher-order representation.
\end{proof}

\begin{corollary}\label{c:mismatch}
The mismatch case raises the same difficulty.
\end{corollary}

\begin{proof}
Suppose $H$ represents $S$ as seeing red while $S$ is seeing green. By Principle~\ref{p:same}, it is for $S$ as if she sees red. The state of seeing red does not exist; the state of seeing green exists but is not represented as it is. If the conscious state is the green state, it is conscious with a character it does not have, so its character as conscious is fixed by $H$ and not by it, and the green state contributes nothing that the targetless case lacks. If the conscious state is the red state, it does not exist. Either way the case reduces to Lemma~\ref{l:three}.
\end{proof}

\section{Objections and replies}

\begin{objection}[Rosenthal]\label{o:ros}
The theory never said that the first-order state must exist. What it is like for the subject is how she represents herself, and a targetless thought makes it seem to her that she is in a state; that seeming is all that consciousness of the state ever was \citep{Rosenthal2011}.
\end{objection}

\begin{reply}
This is option (b), stated as a theory of seeming. If \enquote{the conscious state} is the state the subject seems to be in, then in the targetless case the conscious state is non-existent, and the theory explains phenomenal character by states that need not exist. If instead what is phenomenally present is the seeming, the seeming is $H$, and we are in option (c).
\end{reply}

\begin{objection}[Weisberg]\label{o:weis}
\enquote{Conscious state} is a technical notion; the theory may say that the subject is conscious, in the targetless case, without there being a conscious state \citep{Weisberg2011}.
\end{objection}

\begin{reply}
The proposal grants that there is something it is like for the subject and denies that any state of hers is phenomenally conscious. By Principle~\ref{p:exist}, if there is something it is like for her, some existing state of hers makes it so; the only candidate is $H$, and we are again in option (c).
\end{reply}

\begin{objection}[Brown]\label{o:brown}
Restate the theory: the phenomenally conscious state is the higher-order representation itself, which need not be represented by a further state \citep{Brown2015}.
\end{objection}

\begin{reply}
This is option (c) embraced. It abandons Definition~\ref{d:ho}: what makes the conscious state conscious is no longer that it is represented. The resulting view may be defensible, but it is not a higher-order theory in the sense at issue, and it owes an account of why a representation of a first-order state should be phenomenally conscious when the first-order state need not be.
\end{reply}

\begin{objection}[Rarity]\label{o:rare}
Misrepresentation is rare and abnormal; a theory of normal consciousness need not handle it.
\end{objection}

\begin{reply}
The theories claim to say what consciousness is, not what it normally accompanies, and a constitutive claim must hold in the possible cases. Assumption~\ref{a:poss} is granted by the theories, and their own empirical motivation rests on dissociations of first-order processing and awareness.
\end{reply}

\section{Conclusion}

Higher-order theories make consciousness a matter of how a subject represents herself. When that representation is empty, they must either deny what they say about subjective character, locate consciousness in the representing state and so give up their defining claim, or posit conscious states that do not exist. The last is the only option consistent with the theories, and it is unacceptable. Consciousness is therefore not a higher-order representation.

\begin{thebibliography}{99}\small
\bibitem[Armstrong(1968)]{Armstrong1968} Armstrong, D.~M. (1968). \emph{A Materialist Theory of the Mind}. Routledge \& Kegan Paul.
\bibitem[Block(2011)]{Block2011} Block, N. (2011). The higher order approach to consciousness is defunct. \emph{Analysis}, 71(3), 419--431.
\bibitem[Brown(2015)]{Brown2015} Brown, R. (2015). The HOROR theory of phenomenal consciousness. \emph{Philosophical Studies}, 172(7), 1783--1794.
\bibitem[Byrne(1997)]{Byrne1997} Byrne, A. (1997). Some like it HOT: Consciousness and higher-order thoughts. \emph{Philosophical Studies}, 86(2), 103--129.
\bibitem[Carruthers(2000)]{Carruthers2000} Carruthers, P. (2000). \emph{Phenomenal Consciousness: A Naturalistic Theory}. Cambridge University Press.
\bibitem[Lau and Rosenthal(2011)]{LauRosenthal2011} Lau, H. and Rosenthal, D. (2011). Empirical support for higher-order theories of conscious awareness. \emph{Trends in Cognitive Sciences}, 15(8), 365--373.
\bibitem[Lycan(1996)]{Lycan1996} Lycan, W.~G. (1996). \emph{Consciousness and Experience}. MIT Press.
\bibitem[Neander(1998)]{Neander1998} Neander, K. (1998). The division of phenomenal labor: A problem for representational theories of consciousness. \emph{Philosophical Perspectives}, 12, 411--434.
\bibitem[Rosenthal(2005)]{Rosenthal2005} Rosenthal, D.~M. (2005). \emph{Consciousness and Mind}. Oxford University Press.
\bibitem[Rosenthal(2011)]{Rosenthal2011} Rosenthal, D. (2011). Exaggerated reports: Reply to Block. \emph{Analysis}, 71(3), 431--437.
\bibitem[Weisberg(2011)]{Weisberg2011} Weisberg, J. (2011). Misrepresenting consciousness. \emph{Philosophical Studies}, 154(3), 409--433.
\end{thebibliography}
\end{document}
